A metodologia de pesquisa adotada neste trabalho foi divida em etapas distintas que possibilitam uma análise detalhada do impacto da substituição de Online Codes por RaptorQ Codes na integridade, disponibilidade e eficiência de codificação e decodificação dos dados.

Antes de realizar a proposta de substiuição, foi realizada uma revisão da literatura sobre Online Codes e RaptorQ Codes, bem como uma análise das soluções existentes que utilizam esses códigos para garantir a integridade e disponibilidade dos dados. O objetivo é compreender o estado da arte e identificar lacunas que justificam a substituição de Online Codes por RaptorQ Codes.

As bases de dados utilizadas para a revisão da literatura incluem Scopus, Web of Science, Google Scholar e IEEE Xplore. As palavras-chave utilizadas na busca incluem "Online Codes", "RaptorQ Codes", "proof of retrievability", "erasure codes", "data integrity" e "data availability", bem como seus sinônimos e variações em português e inglês. A busca foi realizada de forma sistemática, considerando artigos publicados nos últimos 8 anos para garantir a relevância e atualidade das informações.

Após a revisão da literatura e a análise das soluções existentes, foi possível identificar lacunas e oportunidades para a substituição de Online Codes por RaptorQ Codes, justificando a proposta apresentada neste trabalho.

\section{Provas de Integridade e Disponibilidade}
O que fundamenta HySAIL são suas provas de integridade e disponibilidade dos dados, que garante que os dados armazenados permanecem corretos e acessíveis mesmo diante de falhas ou ataques maliciosos. Contudo, a prova considera um cenário onde Online Codes são utilizados, sendo necessário avaliar como a substituição por RaptorQ Codes impacta essas garantias. Essa avaliação envolverá a análise teórica das propriedades de RaptorQ Codes em relação à integridade e disponibilidade, bem como experimentos que simulem falhas e ataques para verificar se as garantias são mantidas.

\section{Avaliação de RaptorQ Codes}
Considerando o cenário onde RaptorQ Codes pode substituir Online Codes, é necessário avaliar o impacto dessa substituição na eficiência de codificação e decodificação, bem como na robustez frente a perdas de pacotes. Serão realizados experimentos que medem o overhead de recepção, o tempo de decodificação e a taxa de recuperação de símbolos de origem, permitindo uma análise comparativa entre as duas abordagens.
